% !TEX program = xelatex
% 中文 LNCS/ICONIP-style paper. Layout follows the referenced main.tex.
\documentclass[runningheads]{llncs}

\usepackage{fontspec}
\usepackage{xeCJK}
\usepackage{graphicx}
\usepackage{booktabs}
\usepackage{multirow}
\usepackage{array}
\usepackage{amsmath}
\usepackage{amssymb}
\usepackage{url}

\setCJKmainfont[Path=/System/Library/Fonts/Supplemental/]{Songti.ttc}
\setCJKsansfont[Path=/System/Library/Fonts/Supplemental/]{Songti.ttc}
\setCJKmonofont[Path=/System/Library/Fonts/Supplemental/]{Songti.ttc}

\urlstyle{rm}

\title{面向零样本异常检测的小波监督测试时原型适配}
\titlerunning{小波监督测试时原型适配}
\author{Anonymous Author(s)}
\authorrunning{Anonymous Author(s)}
\institute{Anonymous Institution}

\newcommand{\simi}{\operatorname{sim}}
\newcommand{\normop}{\operatorname{norm}}

\begin{document}
\maketitle

\begin{abstract}
零样本工业异常检测需要在不使用目标类别训练图像和异常标注的条件下，同时完成图像级异常判别与像素级异常定位。现有 CLIP-based 方法通常依赖固定 normal/abnormal 文本原型，但固定原型难以覆盖每张测试图像中的实例化缺陷外观，也容易把真实缺陷、正常材料纹理和结构边界混淆。本文提出 Wavelet-Supervised Test-Time Prototype Adaptation (WPTA)，将 Haar 小波线索用作测试时原型适配的 patch evidence reliability，而不是作为直接叠加到最终异常图上的频域分数。WPTA 冻结 CLIP，在 CLIP patch feature grid 上构造 boundary-aware wavelet reliability，并与初始语义异常先验共同选择 visual normal/abnormal anchors，再通过 confidence-gated conservative update 校准原型。MVTec 和 VisA 上的受控消融表明，direct wavelet fusion 会损害定位表现，而 semantic-spectral evidence selection 与 conservative calibration 能改善测试时原型适配。五个工业基准上，最终校准推理系统相对当前固定 AnomalyCLIP baseline 平均提升 $+1.0/+4.5/+3.4/+2.8$，指标顺序为 P-AUROC/P-AUPRO/I-AUROC/I-AP。
\keywords{零样本异常检测 \and CLIP \and 测试时适配 \and 小波可靠性}
\end{abstract}

\section{引言}

零样本异常检测 (ZSAD) 旨在不使用目标类别训练图像的条件下，对未见类别进行异常判别和异常定位。CLIP-based ZSAD 方法通过比较视觉特征与 normal/abnormal 文本原型，使这一设置具备了可行性。然而，固定 abnormal prototype 只能提供通用异常方向，不能精确匹配当前图像中的具体缺陷外观，例如细划痕、缺口、污染或局部纹理破坏。因此，真实缺陷可能被低估，正常边界、反光或材料纹理也可能被误激活。

频域和小波线索提供了有用但危险的补充信息。许多工业缺陷表现为局部高频扰动，Haar 小波可以在特征网格上分离高频细节与低频结构。然而，高频响应并不等价于异常证据。普通物体轮廓、反光表面和细粒度正常纹理同样会产生强小波响应。受控实验支持这一点：direct wavelet fusion 使 MVTec 的 pixel-level (AUROC, AUPRO) 从 $(91.2,83.2)$ 降到 $(88.7,80.4)$，也使 VisA 从 $(95.5,86.7)$ 降到 $(94.6,85.1)$。

本文因此把小波线索从 final-map scoring 重新定位为测试时原型适配的 evidence reliability。WPTA 首先用固定 normal/abnormal prototypes 得到初始语义异常先验 $S_0$；随后对 CLIP patch feature grid 做 Haar 分解，估计高频纹理能量与低频结构边界，并构造 boundary-aware reliability map。语义先验 $S_0$ 提供 anomaly direction，小波可靠性 $W$ 约束 local evidence 是否可信。最终，WPTA 聚合被两者共同支持的 visual anchors，并在不反向传播、不更新 CLIP 参数的条件下保守校准 prototypes。

本文将证据分为两层。第一层是 MVTec/VisA 受控消融，用于验证 WPTA 机制：semantic-only prototype adaptation 强于固定原型，wavelet-guided reliability 进一步改善定位，conservative update 带来额外收益。第二层是五个工业基准上的最终校准推理系统结果。由于五数据集最终系统使用了数据集特定校准模块，本文显式区分系统级收益和 WPTA 机制的因果证据。

本文贡献如下：
\begin{itemize}
  \item 提出 WPTA，将 Haar 小波线索转化为测试时原型适配的 patch evidence reliability，而不是直接进行 final-map fusion。
  \item 设计 semantic-spectral evidence selection，用语义异常先验和 boundary-aware wavelet reliability 共同构造 visual normal/abnormal anchors。
  \item 在冻结 CLIP、无目标训练数据、无反向传播的条件下进行 conservative prototype calibration，并通过 MVTec/VisA 受控消融验证机制有效性。
  \item 报告五个工业基准结果，并通过 configuration audit 明确区分校准推理系统收益和 WPTA 机制证据。
\end{itemize}

\begin{figure}[t]
  \centering
  \includegraphics[width=0.92\textwidth]{figures/figure1_motivated_example_mvtec_cable_vector_layout.pdf}
  \caption{动机示例。小波线索更适合作为 prototype adaptation 的 evidence reliability，而不是直接作为 anomaly score。}
  \label{fig:motivation}
\end{figure}

\section{相关工作}

\paragraph{CLIP-based zero-shot anomaly detection.}\mbox{}\\
\noindent
CLIP~\cite{radford2021clip} 为开放词汇识别提供了图文对齐基础。WinCLIP~\cite{jeong2023winclip} 使用 window-level CLIP matching 进行 zero-/few-shot anomaly classification and segmentation。AnomalyCLIP~\cite{zhou2024anomalyclip} 学习 object-agnostic normal/abnormal prompts，AdaCLIP~\cite{cao2024adaclip} 通过 hybrid learnable prompts 适配 CLIP。CLIP-AD~\cite{chen2024clipad} 从 language-guided staged dual-path modeling 角度改善零样本异常检测。这些方法证明了 CLIP semantic prototypes 的价值，但大多使用固定或训练得到的 prototypes，缺少来自当前测试图像局部证据的 training-free calibration。

\paragraph{测试时适配。}\mbox{}\\
\noindent
Test-time adaptation 利用测试样本自身信息缓解分布偏移。TPT~\cite{shu2022tpt} 通过 confidence selection 和 view consistency 在测试时调整 vision-language model prompts。异常检测中的测试时适配更困难，因为异常区域稀疏、未知，并常与正常边界或纹理相邻。WPTA 不更新 CLIP 参数，也不从整张图像直接适配；它只聚合同时被语义先验和小波可靠性支持的 patch evidence。

\paragraph{频域与小波线索。}\mbox{}\\
\noindent
小波表示提供了低频结构和高频细节的多分辨率分解~\cite{daubechies1988orthonormal,mallat1989theory}。在工业异常定位中，高频扰动可能对应划痕、裂纹或纹理破坏；但直接频域评分容易受到正常边界和正常纹理干扰。WPTA 将小波线索限制为 evidence reliability，并让校准后的 CLIP prototypes 负责最终 anomaly decision。

\section{方法}

\subsection{问题设置与整体框架}

给定测试图像 $x$，ZSAD 需要预测像素级 anomaly map $M$ 和图像级 anomaly score $s_{\mathrm{img}}$。冻结的 CLIP image encoder 提取 patch features $\{f_i\}_{i=1}^{N}$，text encoder 提供 normal prototype $t_n$ 和 abnormal prototype $t_a$。WPTA 在推理阶段不使用目标域训练图像、异常标注、反向传播或 CLIP 参数更新。

如图~\ref{fig:architecture} 所示，WPTA 包含四个步骤。第一，initial semantic scoring 计算 patch-level anomaly prior $S_0$。第二，boundary-aware Haar wavelet reliability 在 CLIP patch feature grid 上估计 reliability map $W$。第三，semantic-spectral evidence selection 聚合 visual normal/abnormal anchors。第四，conservative prototype calibration 用校准后的 prototypes 重新计算 anomaly probabilities。后续小节按这一执行顺序展开，并分别说明每个模块解决的误差来源。

\begin{figure}[t]
  \centering
  \includegraphics[width=0.92\textwidth]{figures/figure2_wpta_method_overview.pdf}
  \caption{WPTA 方法总览。小波线索用于监督测试时原型适配的 evidence reliability，而不是直接生成最终 anomaly map。}
  \label{fig:architecture}
\end{figure}

\subsection{初始语义异常分数}

初始语义异常分数用于提供 evidence selection 的异常方向。对每个 patch feature $f_i$，WPTA 通过 normal/abnormal prototypes 的相似度 softmax 计算初始语义异常先验：
\begin{equation}
S_0(i)=
\frac{\exp(\simi(f_i,t_a)/\tau)}
{\exp(\simi(f_i,t_a)/\tau)+\exp(\simi(f_i,t_n)/\tau)} .
\end{equation}
其中 $\simi(\cdot,\cdot)$ 表示归一化相似度，$\tau$ 为温度系数。$S_0$ 只提供语义方向，而不直接作为最终校准分数。这样做避免了一个常见失败模式：仅根据固定 abnormal prototype 的响应选择 patch，会把正常结构边界或材料纹理误当作异常 evidence。

\subsection{Prototype mismatch 与适配目标}

原型适配的必要性来自固定文本原型和当前测试图像视觉证据之间的 mismatch。固定 prototypes 实际上用如下 margin 判断 patch 是否更接近异常语义：
\begin{equation}
m_0(i)=\simi(f_i,t_a)-\simi(f_i,t_n).
\end{equation}
当通用 abnormal prototype $t_a$ 与当前图像中的具体缺陷外观不匹配时，真实异常 patch 的 $m_0(i)$ 可能偏低；当 normal prototype $t_n$ 无法覆盖当前材料纹理或结构边界时，正常 patch 的 $m_0(i)$ 也可能被错误抬高。因此，仅依赖固定 prototypes 会同时带来 false negative 和 false positive。

WPTA 将被筛选出的 visual anchors 视为当前测试图像的局部视觉参考。理想的校准 prototypes 应当让 abnormal anchor 更接近 abnormal prototype，同时让 normal anchor 更接近 normal prototype：
\begin{equation}
m'_a=\simi(v_a,t'_a)-\simi(v_a,t'_n),\qquad
m'_n=\simi(v_n,t'_n)-\simi(v_n,t'_a).
\end{equation}
也就是说，校准目标不是任意移动 prototypes，而是在保持文本先验的同时减小 anchor-prototype mismatch。一个等价的约束目标可以写为
\begin{equation}
\mathcal{J}(t'_a,t'_n)=
\|t'_a-v_a\|_2^2+\|t'_n-v_n\|_2^2+
\lambda_a\|t'_a-t_a\|_2^2+
\lambda_n\|t'_n-t_n\|_2^2 .
\end{equation}
前两项把 prototypes 拉向当前图像的 visual anchors，后两项把 prototypes 约束在原始文本先验附近，防止无标签测试时出现 prototype drift。WPTA 使用的 confidence-gated interpolation 可以看作该目标的一步保守近似：当 evidence 可靠时允许小幅校准，当 evidence 不可靠时保留原始文本 prototypes。

\subsection{Boundary-aware Haar wavelet reliability}

Boundary-aware Haar wavelet reliability 用于判断局部 evidence 是否可信，而不是产生新的异常分数。WPTA 将 patch features reshape 为 $F\in\mathbb{R}^{H\times W\times C}$，并进行一级 Haar DWT，得到 $LL$、$LH$、$HL$ 和 $HH$ 分量。实现中，各分量响应被映射回原 CLIP patch grid，以便与 $S_0$ 和 patch features 对齐。高频纹理能量定义为
\begin{equation}
HF(i)=\frac{1}{C}\sum_{c=1}^{C}\left(|LH_i^c|+|HL_i^c|+|HH_i^c|\right).
\end{equation}
由于高频能量也可能来自正常结构边界，WPTA 从 $LL$ 估计低频结构边缘响应 $LF_{\mathrm{edge}}$，并构造
\begin{equation}
W(i)=\normop(HF(i))\cdot\left(1-\normop(LF_{\mathrm{edge}}(i))\right).
\end{equation}
因此，高频纹理扰动会提高 evidence reliability，而低频结构边界会抑制 pseudo evidence。与 direct wavelet fusion 不同，$W$ 不直接与 anomaly map 相加，而只参与后续 evidence selection。

\subsection{Semantic-spectral evidence selection}

Semantic-spectral evidence selection 将语义方向和局部可靠性同时作为筛选条件。Abnormal evidence 应当语义上接近 abnormal prototype，并由可靠局部扰动支持。Normal evidence 应当语义上接近 normal prototype，且不被异常纹理或结构边界污染。概念上，evidence weights 写为
\begin{equation}
q_a(i)=S_0(i)^\gamma\cdot\rho(W(i)),\qquad
q_n(i)=(1-S_0(i))^\gamma\cdot\rho(1-W(i)).
\end{equation}
其中 $\rho(\cdot)$ 表示单调可靠性调制，$\gamma$ 控制语义先验的锐化程度。实现中使用 top-$k$ selection、confidence thresholds、normalization 和 confidence gates 控制 evidence 质量；默认 evidence top-$k$ ratio 为 $0.2$。这一设计的关键点是：语义分数决定 patch 是否像异常，小波可靠性决定该 patch 是否值得用于校准 prototypes。

\subsection{Visual anchors 与保守校准}

Visual anchors 将稀疏 patch evidence 压缩为当前测试图像的原型校准信号。被选择的 evidence patches 聚合为 visual anchors：
\begin{equation}
v_a=\frac{\sum_i q_a(i)f_i}{\sum_i q_a(i)},\qquad
v_n=\frac{\sum_i q_n(i)f_i}{\sum_i q_n(i)}.
\end{equation}
随后，prototypes 被保守校准：
\begin{equation}
t'_a=\operatorname{Norm}\left((1-\alpha)t_a+\alpha v_a\right),\qquad
t'_n=\operatorname{Norm}\left((1-\beta)t_n+\beta v_n\right).
\end{equation}
更新强度由 confidence gates 控制。上式给出通用形式；在默认受控设置中，$\alpha_0=0.0$、$\beta_0=0.01$，因此 abnormal evidence 主要控制更新门控，实际 prototype displacement 主要作用于 normal prototype。这样可以降低无标签测试时校准中的 prototype drift 风险。

校准后，WPTA 用 $t'_a$ 和 $t'_n$ 重新计算 patch-level anomaly probability：
\begin{equation}
S(i)=
\frac{\exp(\simi(f_i,t'_a)/\tau)}
{\exp(\simi(f_i,t'_a)/\tau)+\exp(\simi(f_i,t'_n)/\tau)} .
\end{equation}
随后将 $S(i)$ reshape 并上采样为图像分辨率的 anomaly map。若启用 pixel-to-image fusion，image-level score 会结合 CLIP image-text abnormal probability 与 anomaly map 的 top-$k$ pixel score；若不启用，则 image-level score 直接来自 image-text probability。

\section{实验}

\subsection{实验设置}

\subsubsection{数据集与指标}
\mbox{}\\
\noindent
本文在五个工业异常检测基准上评估：MVTec AD~\cite{bergmann2019mvtec,bergmann2021mvtec}、VisA~\cite{zou2022visa}、MPDD~\cite{jezek2021deep}、BTAD~\cite{mishra2021vt} 和 DTD-Synthetic~\cite{aota2023zero}。MVTec 和 VisA 用于 WPTA 受控机制消融；五数据集表报告最终校准推理系统的系统级结果。

本文报告 Pixel-AUROC、Pixel-AUPRO、Image-AUROC 和 Image-AP。Pixel-AUPRO 更直接反映区域级异常定位质量，因此是主要分析指标。所有 slash-form 数值的顺序固定为 P-AUROC / P-AUPRO / I-AUROC / I-AP。所有结果均为确定性报告，不包含多 seed 置信区间或统计显著性检验。

\subsubsection{实现细节}
\mbox{}\\
\noindent
所有实验使用 AnomalyCLIP checkpoint \texttt{9\_12\_4\_multiscale/epoch\_15.pth} 和 CLIP ViT-L/14@336px。图像 resize 并 center-crop 到 $518\times518$，使用 OpenAI CLIP normalization。缓存特征来自 layers $\{6,12,18,24\}$，受控 WPTA 消融聚合 feature map layers $\{1,2,3\}$。Layer fusion 使用 summation，Gaussian smoothing 使用 $\sigma=5$，AUPRO 使用 200 个 thresholds。

Prototype adaptation 使用 temperature $0.07$、semantic sharpening $\gamma=1.0$、wavelet evidence exponent $\eta=1.0$ 和 evidence top-$k$ ratio $0.2$。Boundary-aware wavelet reliability 使用 boundary-aware mode，mix weight 为 $0.05$。Conservative update 的默认强度为 $\alpha_0=0.0$、$\beta_0=0.01$。除固定 baseline 行外，受控消融中的 prototype/fusion variants 使用相同 multi-crop 和 pixel-to-image 设置：multi-crop weight 为 $0.50$，pixel-to-image weight 为 $0.10$，pixel-to-image top-$k$ ratio 为 $0.01$。

\subsection{结果}

\subsubsection{工业主结果}
\mbox{}\\
\noindent
表~\ref{tab:industry_result} 遵循参考 LNCS 论文的紧凑多方法布局。它同时报告代表性 CLIP-based ZSAD 方法、当前固定 AnomalyCLIP baseline 和 WPTA 结果。外部方法用于 protocol-reference positioning；同协议主 claim 仍来自当前固定 AnomalyCLIP baseline 的提升和后续受控消融。由于外部方法与当前结果可能存在训练条件、split、输入尺寸和后处理差异，表中的 best/second-best 标记仅用于读表定位，不作为严格排序结论。

\begin{table}[htbp]
  \centering
  \caption{五个工业基准上的 zero-shot anomaly detection 结果。每个单元格报告一组指标。收集数值中的 best 和 second-best 分别用 \textbf{bold} 与 \underline{underline} 标记。$^{\dagger}$ 表示参考表中复现结果，$^{\ddagger}$ 表示本文当前固定 baseline 评估。}
  \label{tab:industry_result}
  \scriptsize
  \resizebox{\textwidth}{!}{%
    \begin{tabular}{llllllllll}
      \toprule
      \multicolumn{10}{c}{\textbf{Image-level results (Image-AUROC, Image-AP)}} \\
      \midrule
      Dataset & CLIP & WinCLIP & VAND & CoOp & AdaCLIP & AnomalyCLIP$^{\ddagger}$ & AA-CLIP$^{\dagger}$ & TAAP/INPC & WPTA \\
      \midrule
      MVTec AD & (74.1, 87.6) & (91.8, 96.5) & (86.1, 93.5) & (88.8, 94.8) & (89.2, 96.4) & (91.6, 96.4) & (90.5, 94.9) & (\underline{92.8}, \underline{96.7}) & (\textbf{94.5}, \textbf{97.6}) \\
      VisA & (66.4, 71.5) & (78.1, 81.2) & (78.0, 81.4) & (62.8, 68.1) & (\textbf{85.8}, 84.9) & (82.0, 85.3) & (\underline{84.6}, 82.2) & (83.3, \underline{85.9}) & (\underline{84.6}, \textbf{87.4}) \\
      MPDD & (54.3, 65.4) & (63.6, 69.9) & (73.0, 80.2) & (55.1, 64.2) & (76.0, 80.4) & (73.7, 76.5) & (75.1, 80.1) & (\textbf{81.8}, \textbf{84.8}) & (\underline{77.8}, \underline{82.3}) \\
      BTAD & (34.5, 52.5) & (68.2, 70.9) & (73.6, 68.6) & (66.8, 77.4) & (88.6, 92.4) & (89.1, 91.0) & (\textbf{94.8}, \textbf{97.9}) & (93.3, 93.1) & (\underline{93.9}, \underline{94.9}) \\
      DTD-Synthetic & (71.6, 85.7) & (93.2, 92.6) & (86.4, 95.0) & (-, -) & (95.5, 97.0) & (94.5, 97.7) & (93.3, 97.8) & (\textbf{97.5}, \textbf{99.0}) & (\underline{96.9}, \underline{98.7}) \\
      \midrule
      \textbf{Average} & (60.2, 72.5) & (79.0, 82.2) & (79.4, 83.7) & (68.4, 76.1) & (87.0, 90.2) & (86.2, 89.4) & (87.7, 90.6) & (\textbf{89.7}, \underline{91.9}) & (\underline{89.5}, \textbf{92.2}) \\
      \midrule
      \multicolumn{10}{c}{\textbf{Pixel-level results (Pixel-AUROC, Pixel-AUPRO)}} \\
      \midrule
      MVTec AD & (38.4, 11.3) & (85.1, 64.6) & (87.6, 44.0) & (33.3, 6.7) & (88.7, 37.8) & (91.2, 83.2) & (\textbf{91.9}, 84.6) & (91.5, \underline{85.5}) & (\underline{91.8}, \textbf{85.6}) \\
      VisA & (46.6, 14.8) & (79.6, 56.8) & (94.2, 86.8) & (24.2, 3.8) & (95.5, 56.8) & (95.5, 86.7) & (95.5, 83.0) & (\underline{95.6}, \underline{88.3}) & (\textbf{96.2}, \textbf{91.3}) \\
      MPDD & (62.1, 33.0) & (76.4, 48.9) & (94.1, 83.2) & (15.4, 2.3) & (96.1, 60.3) & (\underline{96.9}, 84.6) & (96.7, 76.5) & (96.7, \textbf{89.9}) & (\textbf{97.3}, \textbf{89.9}) \\
      BTAD & (30.6, 4.4) & (72.7, 27.3) & (60.8, 25.0) & (28.6, 3.8) & (92.1, 32.5) & (93.5, 70.5) & (\textbf{97.0}, 69.0) & (94.1, \underline{77.1}) & (\underline{96.3}, \textbf{78.2}) \\
      DTD-Synthetic & (33.9, 12.5) & (83.9, 57.8) & (95.3, 86.9) & (-, -) & (97.7, 75.0) & (97.4, 89.1) & (96.4, 85.9) & (97.4, \underline{91.5}) & (\textbf{97.9}, \textbf{91.8}) \\
      \midrule
      \textbf{Average} & (42.3, 15.2) & (79.5, 51.1) & (86.4, 65.2) & (25.4, 4.2) & (94.0, 52.5) & (94.9, 82.8) & (\underline{95.5}, 79.8) & (95.1, \underline{86.5}) & (\textbf{95.9}, \textbf{87.4}) \\
      \bottomrule
    \end{tabular}%
  }
\end{table}

从当前固定 AnomalyCLIP baseline 到 WPTA，五数据集平均结果由 $94.9/82.8/86.2/89.4$ 提升到 $95.9/87.4/89.5/92.2$，对应 $+1.0/+4.5/+3.4/+2.8$。最大收益来自 Pixel-AUPRO，说明校准推理系统主要改善异常区域覆盖质量。MPDD 和 BTAD 的最终系统未启用 wavelet reliability 或 prototype rectification，因此这些数据集的提升只作为系统级结果解释；WPTA 机制证据由下方 MVTec/VisA 受控消融给出。

\subsubsection{核心组件消融}
\mbox{}\\
\noindent
表~\ref{tab:core_ablation} 在 MVTec 和 VisA 上验证 WPTA 机制。Direct wavelet fusion 低于固定 baseline，说明小波响应不应直接作为最终 anomaly score。Semantic prototype adaptation 已经强于固定原型，wavelet-guided adaptation 与 conservative update 进一步提升 Pixel-AUPRO。

\begin{table}[htbp]
  \centering
  \caption{MVTec 和 VisA 上的核心组件消融。每个单元格报告 (P-AUROC, P-AUPRO) / (I-AUROC, I-AP)。}
  \label{tab:core_ablation}
  \resizebox{\textwidth}{!}{%
    \begin{tabular}{lll}
      \toprule
      Method & MVTec AD & VisA \\
      \midrule
      Baseline & (91.2, 83.2) / (91.6, 96.4) & (95.5, 86.7) / (82.0, 85.3) \\
      Direct wavelet fusion / no adaptation & (88.7, 80.4) / (92.9, 96.9) & (94.6, 85.1) / (81.6, 84.8) \\
      Semantic prototype adaptation & (91.6, 85.2) / (93.7, 97.1) & (96.0, 90.4) / (83.7, 86.9) \\
      Wavelet prototype adaptation w/o conservative & (91.7, 85.8) / (93.9, 97.2) & (96.1, 91.3) / (84.1, 87.0) \\
      Full WPTA controlled setting & \textbf{(91.8, 86.2) / (94.1, 97.4)} & \textbf{(96.2, 91.7) / (84.3, 87.3)} \\
      \bottomrule
    \end{tabular}%
  }
\end{table}

该消融直接支撑本文的核心设计选择。相比 semantic-only prototype adaptation，Full WPTA 在 MVTec/VisA 上分别带来 $+1.0/+1.3$ 的 P-AUPRO 提升；相比去掉 conservative update 的版本，Full WPTA 进一步提升 $+0.4/+0.4$ P-AUPRO。这说明小波可靠性和保守更新的收益不是来自直接增强 anomaly map，而是来自更可靠的测试时 evidence selection 和更受控的 prototype calibration。

\subsubsection{小波设计消融}
\mbox{}\\
\noindent
表~\ref{tab:wavelet_design} 分析 wavelet reliability 的作用。HF-only reliability 带来有限提升，而 boundary-aware reliability 通过抑制结构边界伪证据进一步改善 Pixel-AUPRO。

\begin{table}[htbp]
  \centering
  \caption{小波可靠性设计消融。每个单元格报告 (P-AUROC, P-AUPRO) / (I-AUROC, I-AP)。}
  \label{tab:wavelet_design}
  \resizebox{\textwidth}{!}{%
    \begin{tabular}{lll}
      \toprule
      Wavelet setting & MVTec AD & VisA \\
      \midrule
      Semantic-only prototype adaptation & (91.6, 85.2) / (93.7, 97.1) & (96.0, 90.4) / (83.7, 86.9) \\
      Direct wavelet fusion & (88.7, 80.4) / (92.9, 96.9) & (94.6, 85.1) / (81.6, 84.8) \\
      HF-only W + prototype adaptation & (91.6, 85.3) / (94.0, 97.2) & (96.0, 90.8) / (84.0, 86.9) \\
      Boundary-aware W + prototype adaptation & (91.7, 85.7) / (93.8, 97.3) & (96.1, 91.2) / (83.9, 87.1) \\
      Full boundary-aware W + conservative & \textbf{(91.8, 86.2) / (94.1, 97.4)} & \textbf{(96.2, 91.7) / (84.3, 87.3)} \\
      \bottomrule
    \end{tabular}%
  }
\end{table}

小波设计消融进一步说明，高频响应有用但不充分。HF-only W 相比 semantic-only 的 P-AUPRO 提升较小；引入 boundary-aware 抑制后，MVTec/VisA 的 P-AUPRO 从 $85.3/90.8$ 提升到 $85.7/91.2$。这与方法设计一致：正常结构边界会产生高频响应，因此需要利用低频结构边缘抑制伪 evidence。

\subsubsection{配置审计}
\mbox{}\\
\noindent
表~\ref{tab:final_config} 记录每个 dataset-specific final system 启用的模块，避免把五数据集所有收益都归因于 WPTA 机制。

\begin{table}[htbp]
  \centering
  \caption{五数据集最终校准推理系统的配置审计。}
  \label{tab:final_config}
  \resizebox{\textwidth}{!}{%
    \begin{tabular}{lccccc}
      \toprule
      Dataset & Result & Wavelet reliability & Prototype rectification & Multi-crop & Pixel-to-image \\
      \midrule
      MVTec AD & 91.8 / 85.6 / 94.5 / 97.6 & yes & yes & yes & yes \\
      VisA & 96.2 / 91.3 / 84.6 / 87.4 & yes & yes & yes & yes \\
      MPDD & 97.3 / 89.9 / 77.8 / 82.3 & no & no & yes & yes \\
      BTAD & 96.3 / 78.2 / 93.9 / 94.9 & no & no & yes & yes \\
      DTD-Synthetic & 97.9 / 91.8 / 96.9 / 98.7 & yes & no & yes & yes \\
      \bottomrule
    \end{tabular}%
  }
\end{table}

\section{讨论}

\paragraph{为什么使用 wavelet reliability 而不是 wavelet scoring。}\mbox{}\\
\noindent
小波提供的是局部频率结构，而不是语义异常判别。Direct-fusion 负对照表明，把 wavelet response 直接加入 anomaly map 会损害定位表现。WPTA 仅将小波用作 evidence reliability，并让校准后的 CLIP prototypes 产生最终 anomaly map。

\paragraph{为什么需要 conservative calibration。}\mbox{}\\
\noindent
无标签测试时校准容易出现 prototype drift。WPTA 使用 confidence-gated conservative updates，使 visual anchors 只有在 evidence 可靠时影响 prototypes。默认设置中 abnormal prototype update strength 为零，abnormal evidence 主要控制 update gate。因此，该方法是保守校准，而不是大幅双向 prototype shift。

\paragraph{范围与局限。}\mbox{}\\
\noindent
五数据集 final systems 使用 dataset-specific configurations，因此它们支持的是相对固定 AnomalyCLIP baseline 的系统级收益，而不是 WPTA 在所有五个数据集上的完全因果隔离。WPTA 机制由 MVTec/VisA 受控消融验证。外部方法比较需要进一步对齐 protocol 后才能用于严格 ranking claim。

\section{结论}

本文提出 WPTA，用于 CLIP-based zero-shot anomaly detection。核心思想是把 Haar 小波线索作为测试时原型适配的 patch evidence reliability，而不是直接进行 final-map fusion。MVTec/VisA 受控消融表明，direct wavelet fusion 会损害定位表现，而 semantic-spectral evidence selection 和 conservative prototype calibration 能改善测试时适配。五个工业基准上，最终校准推理系统相对固定 AnomalyCLIP baseline 获得提升，尤其体现在 Pixel-AUPRO。整体结果支持一个保守结论：小波线索作为 evidence reliability 比作为 anomaly score 更合适，测试时原型校准能够缓解固定 CLIP prototypes 与实例化缺陷外观之间的不匹配。

\bibliographystyle{splncs04}
\bibliography{wpta_references_v0.1}

\end{document}
